\chapter{Luồng hoạt động của chương trình}
\label{chap:program-flow}

\section{Luồng xử lý tổng thể}

Hệ thống được tổ chức theo mô hình frontend--backend. Frontend React đảm
nhận việc thu thập lựa chọn của người dùng, gửi yêu cầu, lưu trạng thái giao
diện và trực quan hóa kết quả. Backend Python kiểm tra yêu cầu, xây dựng đồ
thị theo kịch bản giao thông và thực thi thuật toán tìm kiếm. Vì vậy, frontend
không tự tính tuyến đường và cũng không thay đổi kết quả do backend trả về.

Hình~\ref{fig:flowchart} mô tả luồng xử lý chính của hệ thống.

\begin{figure}[h!]
\centering
\begin{tikzpicture}[
  node distance=0.58cm and 0.5cm,
  box/.style={rectangle, draw, rounded corners, text width=7.0cm,
    minimum height=0.62cm, inner sep=3pt, align=center, font=\scriptsize},
  branch/.style={rectangle, draw, rounded corners, text width=5.0cm,
    minimum height=0.75cm, inner sep=3pt, align=center, font=\scriptsize},
  decision/.style={diamond, draw, aspect=2.6, minimum width=3.7cm,
    align=center, font=\scriptsize},
  arrow/.style={-{Latex[length=1.8mm]}, thick}
]
  \node (start) [box] {Khởi động frontend React};
  \node (load) [box, below=of start]
    {Tự động tải node và cạnh từ hai file GeoJSON};
  \node (input) [box, below=of load]
    {Người dùng chọn điểm đầu, điểm đích, kịch bản, tiêu chí và thuật toán};
  \node (request) [box, below=of input]
    {Frontend kiểm tra lựa chọn, tạo JSON và gửi\\
     \texttt{POST /api/routes/solve}};
  \node (connect) [decision, below=0.8cm of request]
    {Kết nối API thành công?};
  \node (backend) [branch, below left=0.8cm and 0.5cm of connect]
    {Backend kiểm tra request, nạp đồ thị theo kịch bản\\
     và gọi thuật toán};
  \node (fixture) [branch, below right=0.8cm and 0.5cm of connect]
    {Chế độ development thử fixture\\
     cho riêng yêu cầu DL01--DL09};
  \node (result) [box, below=2.35cm of connect]
    {Frontend nhận đối tượng \texttt{result} và lưu vào Zustand};
  \node (status) [decision, below=0.75cm of result]
    {\texttt{status} là \texttt{success}?};
  \node (timeline) [branch, below left=0.8cm and 0.5cm of status]
    {Dựng timeline từ \texttt{frontier\_steps}\\
     hoặc \texttt{visited\_order} và phát mô phỏng};
  \node (failure) [branch, below right=0.8cm and 0.5cm of status]
    {Hiển thị \texttt{no\_path}, \texttt{invalid\_input}\\
     hoặc lỗi dịch vụ};
  \node (display) [box, below=2.35cm of status]
    {Hiển thị node, cạnh đang xét, tuyến cuối, chỉ số và giải thích};

  \draw[arrow] (start) -- (load);
  \draw[arrow] (load) -- (input);
  \draw[arrow] (input) -- (request);
  \draw[arrow] (request) -- (connect);
  \draw[arrow] (connect) -- node[above left, font=\scriptsize] {có} (backend);
  \draw[arrow] (connect) -- node[above right, font=\scriptsize] {không} (fixture);
  \draw[arrow] (backend) |- (result);
  \draw[arrow] (fixture) |- (result);
  \draw[arrow] (result) -- (status);
  \draw[arrow] (status) -- node[above left, font=\scriptsize] {có} (timeline);
  \draw[arrow] (status) -- node[above right, font=\scriptsize] {không} (failure);
  \draw[arrow] (timeline) |- (display);
  \draw[arrow] (failure) |- (display);
\end{tikzpicture}
\caption{Luồng xử lý chính của hệ thống tìm đường}
\label{fig:flowchart}
\end{figure}

Khi chạy ở môi trường development, Vite chuyển tiếp các request bắt đầu bằng
\texttt{/api} đến \texttt{http://127.0.0.1:8000}. Nếu API không thể kết nối,
không có \texttt{VITE\_ROUTE\_API\_URL} tùy chỉnh và ứng dụng không phải bản
production, \texttt{routeService.js} mới thử đọc \texttt{data/mock-result.json}.
Fixture này chỉ chấp nhận yêu cầu từ DL01 đến DL09, không có điểm trung gian;
nó phục vụ kiểm thử giao diện và không chứng minh tính tối ưu của thuật toán.

\section{Kiến trúc và các thành phần chính}

Bảng~\ref{tab:program-modules} trình bày các thành phần quan trọng được xác
minh trong source hiện tại. Tên file được giữ nguyên để thuận tiện đối chiếu
với mã nguồn.

\begin{longtable}{>{\raggedright\arraybackslash}p{3.0cm}
                  >{\raggedright\arraybackslash}p{4.6cm}
                  >{\raggedright\arraybackslash}p{7.0cm}}
\caption{Các thành phần chính của chương trình}
\label{tab:program-modules} \\
\toprule
Thành phần & File hoặc thư mục & Trách nhiệm \\
\midrule
\endfirsthead
\toprule
Thành phần & File hoặc thư mục & Trách nhiệm \\
\midrule
\endhead
Ứng dụng React & \path{src/App.jsx} &
Ghép thanh điều khiển, biểu mẫu chọn tuyến, workspace đồ thị, danh sách thuật
toán và các tab kết quả; đồng thời kích hoạt việc nạp dữ liệu đồ thị. \\
Khung giao diện & \path{components/layout/AppShell.jsx} &
Chia giao diện thành vùng điều khiển phía trên, workspace, sidebar và panel
phía dưới. Ba vùng điều khiển, sidebar và panel kết quả có thể thu gọn độc lập. \\
Nhập yêu cầu & \path{RouteSelectionControls.jsx},
\path{routeRequest.js} &
Lưu điểm đầu, điểm đích, điểm trung gian, kịch bản và tiêu chí; kiểm tra hai
đầu mút; sau đó tạo payload dùng ID node thay vì tên hiển thị. \\
Giao tiếp API & \path{useRouteSolver.js},
\path{routeService.js} &
Gửi request đến \texttt{POST /api/routes/solve}, kiểm tra cấu trúc response và
chỉ dùng fixture trong điều kiện development đã nêu ở trên. \\
Trạng thái toàn cục & \path{src/store/useAppStore.js} &
Quản lý dữ liệu đồ thị, thuật toán được chọn, thông số tuyến, kết quả API,
trạng thái request, trạng thái mô phỏng và thời điểm công bố kết quả cuối. \\
Timeline tìm kiếm & \path{graph/searchTimeline.js} &
Chuyển trace backend thành các action trực quan \texttt{expand},
\texttt{consider-edge} và \texttt{frame-complete}; không tính lại tuyến tối ưu. \\
Hiển thị SVG & \path{GraphWorkspace.jsx} và các layer trong
\path{components/graph/} &
Chiếu tọa độ lên SVG, vẽ mạng đường, trạng thái tìm kiếm, node và tuyến cuối;
hỗ trợ hai bố cục Map/Graph, zoom và kéo bản đồ. \\
Các panel thông tin & \path{BottomPanelTabs.jsx} và
\path{components/results/} &
Hiển thị Simulation, Algorithm trace, Output, Breakdown và Human-readable
reasoning. Ba tab kết quả cuối chỉ mở sau khi mô phỏng lần đầu đi đến cuối. \\
API Python & \path{backend/server.py} &
Cung cấp \texttt{GET /api/health}, \texttt{GET /api/algorithms} và
\texttt{POST /api/routes/solve}; kiểm tra payload rồi chuyển yêu cầu cho solver. \\
Bộ điều phối thuật toán & \path{routing/solver.py} &
Chuẩn hóa tên thuật toán/tiêu chí và gọi BFS, DFS, UCS, Dijkstra, A*, Greedy,
Hill Climbing, Nearest Neighbor hoặc Brute Force TSP theo loại yêu cầu. \\
Nạp đồ thị và dữ liệu & \path{routing/graph_loader.py},
\path{data/generated/} &
Đọc CSV, áp dụng điều kiện S0--S4, loại cạnh bị đóng và tạo
\texttt{NetworkX DiGraph}; frontend đọc hai bản GeoJSON tương ứng để vẽ. \\
\bottomrule
\end{longtable}

Backend công bố chín thuật toán qua API, trong đó Greedy Best-First có trong
backend nhưng chưa có thẻ chọn ở \texttt{AlgorithmSidebar.jsx}. Giao diện hiện
cho chọn tám thuật toán: BFS, DFS, UCS, Dijkstra, A*, Nearest Neighbor, Hill
Climbing và Brute Force TSP. Nearest Neighbor và Brute Force TSP là hai chế độ
nhiều địa điểm; Brute Force TSP giới hạn tối đa tám đích.

\section{Luồng dữ liệu giữa frontend và backend}

Luồng dữ liệu của một yêu cầu tìm đường diễn ra như sau:

\begin{enumerate}[itemsep=2pt]
  \item Người dùng chọn các giá trị trên giao diện. Zustand lưu chúng trong
  \texttt{routeSelection} và lưu thuật toán trong \texttt{selectedAlgorithm}.
  \item \texttt{RouteSelectionControls} chỉ cho gửi form khi đã có hai đầu mút
  khác nhau. \texttt{createRouteRequest} chuyển state thành payload gồm
  \texttt{start\_node}, \texttt{goal\_node}, \texttt{visit\_nodes},
  \texttt{algorithm}, \texttt{scenario\_id} và \texttt{optimization}.
  \item Với Nearest Neighbor hoặc Brute Force TSP, frontend hợp nhất các điểm
  trung gian với điểm đích thành \texttt{visit\_nodes}. Với thuật toán một đích,
  trường này được gửi dưới dạng mảng rỗng.
  \item \texttt{routeService} gửi JSON đến API. \texttt{backend/server.py}
  kiểm tra trường bắt buộc và tính hợp lệ của danh sách điểm.
  \item \texttt{routing/solver.py} nạp đồ thị theo kịch bản và tiêu chí,
  sau đó gọi đúng thuật toán. Các thuật toán trả cùng một hợp đồng kết quả.
  \item \texttt{useRouteSolver} chuyển kết quả cho \texttt{useAppStore}. Nếu
  có trace thành công, mô phỏng tự bắt đầu; nếu không có trace, kết quả được
  công bố ngay.
  \item Khi Zustand thay đổi, React render lại các layer SVG, log, tiến độ,
  thông số từng đoạn và phần giải thích mà backend cung cấp.
\end{enumerate}

Đoạn dưới đây là cấu trúc rút gọn từ một response A* thực tế; trong
\texttt{frontier} chỉ giữ lại một phần tử để minh họa cấu trúc của trace.

\begin{verbatim}
{
  "status": "success",
  "algorithm": "A*",
  "scenario_id": "S0",
  "optimization": "balanced",
  "start_node": "DL01",
  "goal_node": "DL09",
  "path_nodes": ["DL01", "DL06", "DL19", "DL09"],
  "path_edges": ["E005", "E041", "E064"],
  "visited_order": ["DL01", "DL02", "DL05", "DL22"],
  "frontier_steps": [{
    "current": "DL01",
    "frontier": [{"node": "DL02", "g_cost": 0.079213,
                  "h_cost": 0.0, "f_cost": 0.079213}],
    "visited": ["DL01"]
  }],
  "metrics": {
    "total_distance_km": 39.429,
    "total_time_min": 68.307,
    "total_cost": 1.34248,
    "explored_nodes": 25
  },
  "segments": [],
  "explanation": "...",
  "optimality_note": "..."
}
\end{verbatim}

Mảng \texttt{segments} trong response đầy đủ chứa dữ liệu từng cạnh; nó chỉ
được rút gọn thành mảng trống trong đoạn minh họa để tránh lặp lại bảng kết quả
ở Chương~\ref{chap:program-instructions}.

\section{Cơ chế mô phỏng quá trình tìm kiếm}

Backend trả toàn bộ kết quả sau khi thuật toán kết thúc; hệ thống không truyền
trực tiếp từng bước trong lúc thuật toán đang chạy. Frontend dùng dữ liệu trace
đã nhận để phát lại quá trình tìm kiếm:

\begin{itemize}[itemsep=3pt]
  \item \texttt{visited\_order} lưu thứ tự các node đã được mở rộng. Nếu không
  có \texttt{frontier\_steps}, frontend dùng danh sách này làm dữ liệu dự phòng.
  \item \texttt{frontier\_steps} chứa node hiện tại, frontier và tập đã thăm ở
  từng frame. Khi entry có \texttt{g\_cost}, \texttt{h\_cost},
  \texttt{f\_cost} hoặc \texttt{priority}, log có thể hiển thị các giá trị đó.
  \item \texttt{path\_nodes} xác định thứ tự node của tuyến cuối, còn
  \texttt{path\_edges} là danh sách ID cạnh có thứ tự dùng để chọn đúng hình học
  từ GeoJSON. Frontend không nối các điểm bằng suy đoán tọa độ.
  \item Action \texttt{expand} đánh dấu bắt đầu mở rộng node hiện tại.
  \texttt{consider-edge} lần lượt làm sáng mỗi cạnh đi ra để minh họa thao tác
  xét/relax cạnh. \texttt{frame-complete} đồng bộ frontier sau khi đã xét xong
  node. Kết quả ADD, UPDATE, KEEP hoặc UNKNOWN chỉ được suy ra từ hai snapshot
  frontier nhằm phục vụ trình bày, không thay thế phép tính backend.
\end{itemize}

Đối với nhánh tìm kiếm dẫn từ điểm bắt đầu đến node đang xét, frontend ưu tiên
\texttt{current\_path\_edges}, \texttt{path\_edges} hoặc quan hệ parent nếu
frame backend có cung cấp. Nếu trace không có các trường này, frontend dựng một
nhánh minh họa từ quan hệ phát hiện node lần đầu trong các snapshot frontier và
chỉ chấp nhận cạnh có hướng tồn tại trong GeoJSON. Nhánh minh họa này giúp người
xem theo dõi quá trình; nó không được dùng làm tuyến kết quả và không tác động
đến \texttt{result.path\_edges}.

Zustand là nguồn state duy nhất của \texttt{simulation.status},
\texttt{simulation.currentStep} và \texttt{simulation.speed}. First action đưa
về action đầu; Previous action lùi một action; Next action tiến một action;
Last action đến action cuối; Play/Pause điều khiển phát tự động; Reset đưa
timeline về đầu nhưng giữ tốc độ đã chọn. Khoảng nghỉ mặc định giữa hai action
là $700/\textit{speed}$ mili giây với các mức $0.5\times$, $1\times$,
$1.5\times$ và $2\times$.

Ba tab Output, Breakdown và Human-readable reasoning bị khóa trước lần đầu
timeline đến cuối. Khi đã công bố, kết quả vẫn còn nếu người dùng quay lại bằng
Previous, First hoặc Reset; nó chỉ bị xóa khi bắt đầu yêu cầu mới, xóa kết quả
hoặc gặp lỗi request mới. Tuyến cuối cũng chỉ hiện khi mô phỏng hoàn tất, ngoại
trừ response thành công không có timeline thì được hiện ngay.

\section{Luồng hiển thị bản đồ và đồ thị}

Hook \texttt{useGraphData} tự động tải 25 node từ
\texttt{nodes\_snapped.geojson} và 114 cạnh có hướng từ
\texttt{edges.geojson} trong thư mục
\path{data/generated/}. Workspace dùng SVG và xếp các
layer theo thứ tự sau:

\begin{enumerate}[itemsep=2pt]
  \item \texttt{RoadNetworkLayer} vẽ mạng đường nền;
  \item \texttt{SearchTraversalLayer} vẽ nhánh tìm kiếm, các cạnh ứng viên và
  cạnh đang được xét;
  \item \texttt{FinalRouteLayer} vẽ các cạnh thuộc \texttt{path\_edges};
  \item \texttt{SearchAnimationLayer} và \texttt{GraphNodeLayer} vẽ trạng thái
  node phía trên các đường.
\end{enumerate}

Chú thích node gồm Unvisited, Frontier, Visited, Current và Final-route node.
Chú thích cạnh gồm Candidate edge, Active relaxation, Normal route segment và
Warning route segment. Một đoạn của tuyến cuối được cảnh báo khi mức ùn tắc từ
4 trở lên hoặc risk lớn hơn 0; bởi vậy tuyến cuối có thể đồng thời có đoạn xanh
và đoạn cam.

Chế độ Map giữ tọa độ địa lý và hình học đường thực. Chế độ Graph dùng bố cục
topology xác định để giãn đều node nhưng vẫn giữ các kết nối có hướng. Người
dùng có thể phóng to/thu nhỏ bằng nút hoặc con lăn chuột trong khoảng từ
$60\%$ đến $500\%$, giữ chuột trái để kéo, nhấn đúp hoặc chọn Reset view để
trở về góc nhìn ban đầu. Các vùng Route controls, Search strategy và Search
details có thể thu gọn độc lập; component vẫn được mount nên lựa chọn, tab,
timeline và trạng thái bản đồ không bị mất.
